% ===============================
% MAT221 -- Real Analysis
% Riemann Integral
% ===============================

\documentclass[12pt,a4paper]{article}

% ===============================
% Metadata
% ===============================
\newcommand*{\CourseCode}{MAT221}
\newcommand*{\CourseTitle}{Real Analysis}
\newcommand*{\Chapter}{Measure Theory Crash Course}
\newcommand*{\Topics}{$\sigma$-Algebras, Measures, Outer Measures, Lebesgue Measure, Non Measurable Sets}
\newcommand*{\DocDate}{September 01, 2026}

% ===============================
% Header
% ===============================
\input{header}

% ===============================
% Document
% ===============================
\begin{document}

\FrontPage
\Large

\begin{heading}
    What is a Measure?
\end{heading}

\clearpage

\begin{heading}
    $\sigma$-Algebras
\end{heading}

\vspace{10pt}

\begin{definition}[$\sigma$-Algebra]
    Let $X$ be a set. A \textbf{$\sigma$-algebra} $\mathcal{A}$ on $X$ is a collection of subsets of $X$ such that:
    \begin{enumerate}
        \item $X \in \mathcal{A}$ (the whole set is in the $\sigma$-algebra),
        \item If $A \in \mathcal{A}$, then $A^c = X \setminus A \in \mathcal{A}$ (closed under complementation),
        \item If $\{A_n\}_{n=1}^{\infty} \subseteq \mathcal{A}$, then $\bigcup_{n=1}^{\infty} A_n \in \mathcal{A}$ (closed under countable unions).
    \end{enumerate}
\end{definition}

\clearpage

\begin{theorem}[Properties of $\sigma$-Algebras]
    Let $\mathcal{A}$ be a $\sigma$-algebra on a set $X$. Then:
    \begin{enumerate}
        \item $\emptyset \in \mathcal{A}$ (the empty set is in the $\sigma$-algebra),
        \item If $A, B \in \mathcal{A}$, then $A \cap B \in \mathcal{A}$ (closed under finite intersections),
        \item If $\{A_n\}_{n=1}^{\infty} \subseteq \mathcal{A}$, then $\bigcap_{n=1}^{\infty} A_n \in \mathcal{A}$ (closed under countable intersections).
    \end{enumerate}
    
\end{theorem}

\clearpage

\begin{theorem}[Intersection of $\sigma$-Algebras]
    Let $\{\mathcal{A}_i\}_{i \in I}$ be a collection of $\sigma$-algebras on a set $X$. Then the intersection $\bigcap_{i \in I} \mathcal{A}_i$ is also a $\sigma$-algebra on $X$.
\end{theorem}

\clearpage

\begin{definition}[Generated $\sigma$-Algebra]
    Let $X$ be a set and let $\mathcal{C}$ be a collection of subsets of $X$. The \textbf{$\sigma$-algebra generated by $\mathcal{C}$}, denoted by $\sigma(\mathcal{C})$, is the smallest $\sigma$-algebra on $X$ that contains all the sets in $\mathcal{C}$. Formally, it is defined as:
    \[
        \sigma(\mathcal{C}) = \bigcap \{\mathcal{A} : \mathcal{A} \text{ is a } \sigma\text{-algebra on } X \text{ and } \mathcal{C} \subseteq \mathcal{A}\}.
    \]
\end{definition}

\clearpage

\begin{definition}[Borel $\sigma$-Algebra]
    The \textbf{Borel $\sigma$-algebra} on the real line $\mathbb{R}$, denoted by $\mathcal{B}(\mathbb{R})$, is the $\sigma$-algebra generated by the open intervals in $\mathbb{R}$. That is,
    \[
        \mathcal{B}(\mathbb{R}) = \sigma\left(\{(a, b) : a < b, a, b \in \mathbb{R}\}\right).
    \]
\end{definition}

\clearpage

\begin{problem}[Size of the Borel $\sigma$-Algebra]
    Show that cardinality of the Borel $\sigma$-algebra $\mathcal{B}(\mathbb{R})$ is $\mathfrak{c}$. In other words, show that $\mathcal{B}(\mathbb{R})$ has the same cardinality as the real numbers.
\end{problem}

\clearpage

\begin{heading}
    Measure Spaces and Measures
\end{heading}

\vspace{10pt}

\begin{definition}[Measurable Space]
    A \textbf{measurable space} is a pair $(X, \mathcal{A})$ where $X$ is a set and $\mathcal{A}$ is a $\sigma$-algebra on $X$. The elements of $\mathcal{A}$ are called \textbf{measurable sets}.
\end{definition}

\clearpage

\begin{definition}[Measure]
    Let $(X, \mathcal{A})$ be a measurable space. A \textbf{measure} $\mu$ on $(X, \mathcal{A})$ is a function $\mu: \mathcal{A} \to [0, \infty]$ such that:
    \begin{enumerate}
        \item $\mu(\emptyset) = 0$ (the measure of the empty set is zero),
        \item (Countable Additivity) For any countable collection of disjoint sets $\{A_n\}_{n=1}^{\infty} \subseteq \mathcal{A}$,
        \[
            \mu\left(\bigcup_{n=1}^{\infty} A_n\right) = \sum_{n=1}^{\infty} \mu(A_n).
        \]
    \end{enumerate}
\end{definition}

\clearpage

\begin{heading}
    Standard Examples of Measures
\end{heading}

\vspace{10pt}

\begin{problem}[Counting Measure]
    Let $X$ be a set and let $\mathcal{A} = \mathcal{P}(X)$ be the power set of $X$. Define the \textbf{counting measure} $\mu: \mathcal{A} \to [0, \infty]$ by
    \[
        \mu(A) = |A| \quad \text{for all } A \in \mathcal{A},
    \]
    where $|A|$ denotes the cardinality of the set $A$. Show that $\mu$ is a measure on $(X, \mathcal{A})$.
\end{problem}

\clearpage

\begin{problem}[Dirac Measure]
    Let $X$ be a set and let $x_0 \in X$. Define the \textbf{Dirac measure} $\delta_{x_0}: \mathcal{P}(X) \to [0, \infty]$ by
    \[
        \delta_{x_0}(A) = 
        \begin{cases}
            1 & \text{if } x_0 \in A, \\
            0 & \text{if } x_0 \notin A,
        \end{cases}
    \]
    for all $A \subseteq X$. Show that $\delta_{x_0}$ is a measure on $(X, \mathcal{P}(X))$.
\end{problem}

\clearpage

\begin{problem}[Borel Measure]
    Let $\mathcal{B}(\mathbb{R})$ be the Borel $\sigma$-algebra on $\mathbb{R}$. Define the \textbf{Borel measure} $\mu: \mathcal{B}(\mathbb{R}) \to [0, \infty]$ by
    \[
        \mu((a, b)) = b - a \quad \text{for all } a < b, a, b \in \mathbb{R},
    \]
    and extend it to all Borel sets using the properties of measures. Show that $\mu$ is a measure on $(\mathbb{R}, \mathcal{B}(\mathbb{R}))$.
\end{problem}

\clearpage

\begin{theorem}[Properties of Measures]
    Let $(X, \mathcal{A}, \mu)$ be a measure space. Then the following properties hold:
    \begin{enumerate}
        \item (Monotonicity) If $A, B \in \mathcal{A}$ and $A \subseteq B$, then $\mu(A) \leq \mu(B)$.
        \item (Subadditivity) For any countable collection of sets $\{A_n\}_{n=1}^{\infty} \subseteq \mathcal{A}$,
        \[
            \mu\left(\bigcup_{n=1}^{\infty} A_n\right) \leq \sum_{n=1}^{\infty} \mu(A_n).
        \]
        \item (Continuity from Below) If $\{A_n\}_{n=1}^{\infty} \subseteq \mathcal{A}$ is an increasing sequence of sets (i.e., $A_n \subseteq A_{n+1}$ for all $n$), then
        \[
            \mu\left(\bigcup_{n=1}^{\infty} A_n\right) = \lim_{n \to \infty} \mu(A_n).
        \]
        \item (Continuity from Above) If $\{A_n\}_{n=1}^{\infty} \subseteq \mathcal{A}$ is a decreasing sequence of sets (i.e., $A_n \supseteq A_{n+1}$ for all $n$) and $\mu(A_1) < \infty$, then
        \[
            \mu\left(\bigcap_{n=1}^{\infty} A_n\right) = \lim_{n \to \infty} \mu(A_n).
        \]
    \end{enumerate}
    
\end{theorem}

\clearpage

\begin{problem}[Borel Measure of $\mathbb{Q}$]
    Show that the Borel measure of the set of rational numbers $\mathbb{Q}$ is zero, i.e., $\mu(\mathbb{Q}) = 0$.
\end{problem}

\clearpage

\begin{problem}[Borel Measure of a Singleton]
    Show that the Borel measure of a singleton set $\{x\}$, where $x \in \mathbb{R}$, is zero, i.e., $\mu(\{x\}) = 0$.
\end{problem}

\clearpage

\begin{problem}[Borel Measure of a countable set]
    Show that the Borel measure of any countable set $A \subseteq \mathbb{R}$ is zero, i.e., $\mu(A) = 0$.
\end{problem}

\clearpage

\begin{heading}
    Outer Measures
\end{heading}

\vspace{10pt}

\begin{definition}[Outer Measure]
    An \textbf{outer measure} on a set $X$ is a function $\mu^*: \mathcal{P}(X) \to [0, \infty]$ such that:
    \begin{enumerate}
        \item $\mu^*(\emptyset) = 0$ (the outer measure of the empty set is zero),
        \item (Monotonicity) If $A, B \subseteq X$ and $A \subseteq B$, then $\mu^*(A) \leq \mu^*(B)$,
        \item (Countable Subadditivity) For any countable collection of sets $\{A_n\}_{n=1}^{\infty} \subseteq X$,
        \[
            \mu^*\left(\bigcup_{n=1}^{\infty} A_n\right) \leq \sum_{n=1}^{\infty} \mu^*(A_n).
        \]
    \end{enumerate}
\end{definition}

\clearpage

\begin{definition}[Length of an Interval in $\mathbb{R}$]
    The \textbf{length} of an interval $I = (a, b)$, where $a < b$, is defined as $\ell(I) = b - a$. For other types of intervals, the length is defined similarly:
    \begin{itemize}
        \item For a closed interval $[a, b]$, $\ell([a, b]) = b - a$.
        \item For a half-open interval $(a, b]$ or $[a, b)$, $\ell((a, b]) = \ell([a, b)) = b - a$.
    \end{itemize}
\end{definition}

\clearpage

\begin{definition}[Volume of a Rectangle in $\mathbb{R}^n$]
    Let $R = \prod_{i=1}^{n} [a_i, b_i]$ be a rectangle in $\mathbb{R}^n$, where $a_i < b_i$ for all $i = 1, 2, \ldots, n$. The \textbf{volume} of the rectangle $R$ is defined as
    \[
        \text{Vol}(R) = \prod_{i=1}^{n} (b_i - a_i).
    \]
\end{definition}

\clearpage

\begin{definition}[Lebesgue outer measure in $\mathbb{R}$]
    The \textbf{Lebesgue outer measure} on $\mathbb{R}$ is the function $\mu^*: \mathcal{P}(\mathbb{R}) \to [0, \infty]$ defined by
    \[
        \mu^*(A) = \inf\left\{\sum_{i=1}^{\infty} \ell(I_i) : A \subseteq \bigcup_{i=1}^{\infty} I_i\right\}
    \]
    for all $A \subseteq \mathbb{R}$, where the infimum is taken over all countable collections of intervals $\{I_i\}_{i=1}^{\infty}$ that cover $A$.
\end{definition}

\clearpage

\begin{definition}[Lebesgue outer measure in $\mathbb{R}^n$]
    The \textbf{Lebesgue outer measure} on $\mathbb{R}^n$ is the function $\mu^*: \mathcal{P}(\mathbb{R}^n) \to [0, \infty]$ defined by
    \[
        \mu^*(A) = \inf\left\{\sum_{i=1}^{\infty} \text{Vol}(R_i) : A \subseteq \bigcup_{i=1}^{\infty} R_i\right\}
    \]
    for all $A \subseteq \mathbb{R}^n$, where the infimum is taken over all countable collections of rectangles $\{R_i\}_{i=1}^{\infty}$ that cover $A$.
\end{definition}

\clearpage

\begin{theorem}[Outer measure of an interval]
    Let $I = (a, b)$ be an interval in $\mathbb{R}$. Then the Lebesgue outer measure of $I$ is equal to its length, i.e.,
    \[
        \mu^*(I) = \ell(I) = b - a.
    \]
\end{theorem}

\vspace{100pt}

\begin{theorem}[Outer measure of a rectangle]
    Let $R = \prod_{i=1}^{n} [a_i, b_i]$ be a rectangle in $\mathbb{R}^n$. Then the Lebesgue outer measure of $R$ is equal to its volume, i.e.,
    \[
        \mu^*(R) = \text{Vol}(R) = \prod_{i=1}^{n} (b_i - a_i).
    \]
\end{theorem}

\clearpage

\begin{theorem}[Translation invariance of outer measure]
    Let $A \subseteq \mathbb{R}$ and $c \in \mathbb{R}$. Then the Lebesgue outer measure of the translated set $A + c = \{x + c : x \in A\}$ is equal to the Lebesgue outer measure of $A$, i.e.,
    \[
        \mu^*(A + c) = \mu^*(A).
    \]
\end{theorem}

\clearpage

\begin{theorem}[Outer measure of a countable set]
    Let $A \subseteq \mathbb{R}$ be a countable set. Then the Lebesgue outer measure of $A$ is zero, i.e.,
    \[
        \mu^*(A) = 0.
    \]
\end{theorem}

\clearpage

\begin{heading}
    Carathéodory's Criterion for Measurability
\end{heading}

\vspace{10pt}

\begin{definition}[Carathéodory Measurable Set]
    A set $E \subseteq \mathbb{R}$ is said to be \textbf{Carathéodory measurable} (or simply \textbf{measurable}) with respect to the Lebesgue outer measure $\mu^*$ if for every set $A \subseteq \mathbb{R}$,
    \[
        \mu^*(A) = \mu^*(A \cap E) + \mu^*(A \cap E^c).
    \]
\end{definition}

\clearpage

\begin{theorem}[Carathéodory's theorem]
    The collection of all Carathéodory measurable sets forms a $\sigma$-algebra, and the restriction of the Lebesgue outer measure $\mu^*$ to this $\sigma$-algebra is a measure, called the \textbf{Lebesgue measure}.
\end{theorem}

\clearpage

\begin{heading}
    Lebesgue Measure
\end{heading}

\vspace{10pt}

\begin{definition}[Lebesgue Measure]
    The \textbf{Lebesgue measure} $\mu$ on $\mathbb{R}$ is the measure defined on the $\sigma$-algebra of Carathéodory measurable sets, which is a subset of the Borel $\sigma$-algebra. For any Carathéodory measurable set $E \subseteq \mathbb{R}$, the Lebesgue measure is given by
    \[
        \mu(E) = \mu^*(E),
    \]
    where $\mu^*$ is the Lebesgue outer measure.
    
\end{definition}

\clearpage

\begin{definition}[Translation Invariance of Lebesgue Measure]
    An important geometric property of Lebesgue measure is its translational in-variance.
    The Lebesgue measure $\mu$ is said to be \textbf{translation invariant} if for any measurable set $E \subseteq \mathbb{R}$ and any real number $c \in \mathbb{R}$, the translated set $E + c = \{x + c : x \in E\}$ is also measurable and
    \[
        \mu(E + c) = \mu(E).
    \]
    
\end{definition}

\clearpage

\begin{definition}[Null Set]
    A set $N \subseteq \mathbb{R}$ is called a \textbf{null set} (or a set of measure zero) if its Lebesgue measure is zero, i.e.,
    \[
        \mu(N) = 0.
    \]
\end{definition}

\clearpage

\begin{definition}[Complete Measure Space]
    A measure space $(X, \mathcal{A}, \mu)$ is said to be \textbf{complete} if every subset of a null set is measurable. That is, if $N \in \mathcal{A}$ with $\mu(N) = 0$ and $S \subseteq N$, then $S \in \mathcal{A}$ and $\mu(S) = 0$.
\end{definition}

\clearpage

\begin{theorem}[Completeness of Lebesgue Measure]
    The Lebesgue measure $\mu$ on $\mathbb{R}$ is a complete measure. That is, if $N \subseteq \mathbb{R}$ is a null set (i.e., $\mu(N) = 0$) and $S \subseteq N$, then $S$ is Lebesgue measurable and $\mu(S) = 0$.
\end{theorem}

\clearpage

\begin{theorem}[Lebesgue Measure of a countable set]
    Let $A \subseteq \mathbb{R}$ be a countable set. Then the Lebesgue measure of $A$ is zero, i.e.,
    \[
        \mu(A) = 0.
    \]
\end{theorem}

\vspace{10pt}

\begin{problem}[Is there any uncountable set of measure zero?]
    Can you give an example of an uncountable set of measure zero?
\end{problem}

\clearpage

\begin{theorem}[Cantor Set definition]
    The Cantor set $C$ is constructed by iteratively removing the open middle third intervals from the closed interval $[0, 1]$. Formally, we define the Cantor set as follows:
    \begin{enumerate}
        \item Start with the closed interval $C_0 = [0, 1]$.
        \item Remove the open middle third interval $(\frac{1}{3}, \frac{2}{3})$ to obtain $C_1 = [0, \frac{1}{3}] \cup [\frac{2}{3}, 1]$.
        \item In each subsequent step, remove the open middle third intervals from each remaining closed interval. For example, in the second step, remove $(\frac{1}{9}, \frac{2}{9})$ and $(\frac{7}{9}, \frac{8}{9})$ from $C_1$ to obtain $C_2 = [0, \frac{1}{9}] \cup [\frac{2}{9}, \frac{1}{3}] \cup [\frac{2}{3}, \frac{7}{9}] \cup [\frac{8}{9}, 1]$, and so on.
        \item The Cantor set is defined as the intersection of all these sets:
        \[
            C = \bigcap_{n=0}^{\infty} C_n.
        \]
    \end{enumerate}
\end{theorem}

\clearpage

\begin{theorem}[Cantor set has measure zero]
    Show that the Cantor set $C$ has Lebesgue measure zero, i.e., $\mu(C) = 0$.
\end{theorem}

\clearpage

\begin{theorem}[Cantor set is uncountable]
    Show that the Cantor set $C$ is uncountable.
\end{theorem}

\clearpage

\begin{problem}[Size of the Lebesgue $\sigma$-Algebra]
    Show that the cardinality of the Lebesgue $\sigma$-algebra is $2^{\mathfrak{c}}$. In other words, show that the Lebesgue $\sigma$-algebra has the same cardinality as the power set of the real numbers.
\end{problem}

\clearpage

\begin{heading}
    Almost Everywhere notation
\end{heading}

\begin{definition}[Almost Everywhere]
    A property $P(x)$ is said to hold \textbf{almost everywhere} (a.e.) on a measure space $(X, \mathcal{A}, \mu)$ if the set of points where $P(x)$ does not hold has measure zero. Formally, $P(x)$ holds almost everywhere if
    \[
        \mu(\{x \in X : P(x) \text{ does not hold}\}) = 0.
    \]
\end{definition}

\clearpage

\begin{heading}
    Is all subsets of $\mathbb{R}^n$ Lebesgue measurable? 
\end{heading}

\vspace{10pt}

\begin{problem}[Non measurable sets]
    Give an example of a non-measurable set in $\mathbb{R}$. (Hint: Consider the Vitali set.)
\end{problem}

\clearpage

\begin{heading}
    Some Important Results and Notes
\end{heading}

\vspace{10pt}

\begin{problem}[Why Countable Additivity is Important?]
    Explain why countable additivity is an important property for a measure. What would go wrong if we allow uncountable additivity instead? Also why finite additivity is not enough to define a measure?
\end{problem}

\clearpage

\begin{problem}[Is all subsets of $\mathbb{R}^n$ Lebesgue measurable?]
    It is not possible to define the Lebesgue measure of all subsets of $\mathbb{R}^n$. Hausdorff (1914) showed that for any dimension $n \geq 1$, there is no countably additive measure defined on all subsets of $\mathbb{R}^n$ that is translation invariant and assigns measure $1$ to the unit cube. In particular, there exist non-measurable sets in $\mathbb{R}^n$.

    He further showed that if $n \geq 3$, there is no such finitely additive measure. This result is dramatized by the Banach-Tarski paradox, which states that a solid ball in 3-dimensional space can be split into a finite number of pieces and reassembled into two solid balls identical to the original. This paradox relies on the existence of non-measurable sets.
\end{problem}

\clearpage

\EndPage

\end{document}